% ============================================================
% Teil B: Schülermaterial -- Durchsage (Vorlesetext) und Regelkarte
% ============================================================
\clearpage
\label{mat:start}
\abschnitt{Durchsage und Hintergrund \textnormal{\normalsize (Vorlesetext für die Lehrkraft)}}
\szene{5.5}

\ifkurz
\subsection*{Jahr 2033 -- Die Oberfläche ist verloren}
Vor elf Jahren veränderte ein globaler Krieg die Welt für immer. Feuerstürme, radioaktive Niederschläge und eine unbekannte Seuche machten die Oberfläche unbewohnbar. Die Überlebenden flohen in die U-Bahn-Schächte und Schutzbunker. Strom gibt es nur stundenweise, Wasser kommt aus alten Filteranlagen, Nahrung wächst in Pilzkellern. Medikamente sind wertvoller als Munition.

Ihr lebt in der Station \textbf{Endstation 7}, einem ehemaligen Verkehrsknotenpunkt mit etwa 120 Menschen. Vor drei Tagen zerstörte eine Explosion die zentrale Pumpenanlage, dann fiel das Stromnetz aus. Die bisherige Stationsleitung ließ Vorräte beschlagnahmen und Proteste niederschlagen. Als sie die Evakuierung ankündigte, war der Verbindungstunnel eingestürzt. Die Leitung ist verschwunden.

Fünf Gemeinschaften kontrollieren die fünf wichtigsten Bereiche der Station. Jede besitzt etwas, das alle zum Überleben brauchen. Keine besitzt alles. Dann knackt der Lautsprecher:
\else
\subsection*{Jahr 2033 -- Die Oberfläche ist verloren}
Vor elf Jahren veränderte ein globaler Krieg die Welt für immer. Feuerstürme, radioaktive Niederschläge und eine unbekannte Seuche machten weite Teile der Oberfläche unbewohnbar. Die wenigen Überlebenden flohen in die U-Bahn-Schächte, Schutzbunker und Versorgungstunnel. Was zunächst als vorübergehende Zuflucht gedacht war, wurde zur neuen Heimat der Menschheit.

Die Metro ist längst keine funktionierende Stadt mehr. Strom gibt es nur stundenweise. Wasser wird durch alte Filteranlagen gepumpt. Nahrung wächst in improvisierten Pilzkellern und Gewächskammern. Medikamente sind wertvoller als Munition. Und draußen in den Tunneln lauern Gefahren, über die niemand gerne spricht.

Ihr lebt in der Station \textbf{Endstation 7}, einem ehemaligen Verkehrsknotenpunkt, der inzwischen etwa 120 Menschen beherbergt. Die Station ist groß genug, um eine Gemeinschaft zu bilden, aber zu klein, um unabhängig zu überleben.

\subsection*{Der Tag, an dem alles zusammenbricht}
Vor drei Tagen ereignete sich eine Katastrophe: Eine Explosion in einem Versorgungstunnel zerstörte die zentrale Pumpenanlage. Kurz darauf fiel das überalterte Stromnetz aus. Die bisherige Stationsleitung, ein kleiner Rat aus ehemaligen Sicherheitskräften und Technikern, versuchte, die Lage unter Kontrolle zu halten. Doch die Vorräte wurden beschlagnahmt, Wasser wurde rationiert und mehrere Menschen wurden bei Protesten verletzt. Als die Stationsleitung die Evakuierung in einen anderen Metroabschnitt ankündigte, stellte sich heraus, dass der Verbindungstunnel eingestürzt war.

Die Leitung ist verschwunden. Niemand weiß, ob sie geflohen ist oder ums Leben kam. Seitdem herrscht Unsicherheit.

Die fünf wichtigsten Bereiche der Station werden von fünf Gemeinschaften kontrolliert. Jede besitzt etwas, das für das Überleben aller unverzichtbar ist. Keine Gemeinschaft verfügt jedoch über alles, was sie benötigt.

Am heutigen Morgen versammelt sich die Bevölkerung auf dem Bahnsteig. Eine alte Notstromanlage liefert noch genug Energie, um eine letzte Durchsage abzuspielen:
\fi

\begin{tcolorbox}[enhanced,colback=black!85,colframe=black,coltext=white,arc=1.5mm,left=5mm,right=5mm,top=3mm,bottom=3mm]
\itshape \enquote{An alle Überlebenden von Endstation~7. Die bisherige Ordnung ist zusammengebrochen. Es gibt keine übergeordnete Verwaltung mehr. Ihr müsst entscheiden, wie ihr künftig zusammenleben wollt. Findet eine gemeinsame politische Ordnung -- oder riskiert, dass die Station auseinanderbricht.}
\end{tcolorbox}

Die Lautsprecher verstummen. \textbf{Ihr habt 48 Stunden, um eine neue Regierung zu bilden.} Wenn ihr euch nicht einigt, reichen die Vorräte nur noch für wenige Tage. Niemand kann garantieren, dass die Station danach noch bewohnbar ist.

% ------------------------------------------------------------
\clearpage
\abschnitt{Regelkarte: So läuft Endstation 7}
\small\setstretch{1.07}

\begin{metrobox}[title={1 \ Marker},fonttitle=\bfseries]{cS}
Jedes Gut wird durch Marker dargestellt: \rN~Nahrung, \rW~Wasser, \rE~Energie, \rM~Medikamente, \rS~Schutzmaterial. Marker dürfen getauscht, verschenkt, verliehen oder gegen politische Zugeständnisse ausgehandelt werden. Alles, was ihr abgebt, bekommt ihr nur zurück, wenn ihr es vereinbart habt.
\end{metrobox}

\begin{metrobox}[title={2 \ Eine Runde (\ifkurz 24 \else 16 \fi Stunden Spielzeit)},fonttitle=\bfseries]{cS}
\phase{1}~\textbf{Produktion:} Die Bank gibt eure Produktionsmarker aus.\\
\phase{2}~\textbf{Handel und Verhandlung:} Nur die Delegation geht zum Konferenztisch. Gespräche zu zweit sind erlaubt.\\
\phase{3}~\textbf{Politische Entscheidungen:} Ab Runde~2 beraten alle Gruppen über gemeinsame Regeln.\\
\phase{4}~\textbf{Betriebskosten:} Jede Gruppe zahlt ihre Betriebskosten an die Bank (steht auf eurer Gruppenkarte).\\
\phase{5}~\textbf{Ereignis und Abrechnung:} Krankheitskarte, Ereignis, Abrechnung der Stabilität.
\end{metrobox}

\begin{metrobox}[title={3 \ Betriebskosten und Stufen},fonttitle=\bfseries]{cE}
\textbf{Stufe~1 Normalbetrieb:} alle Kosten bezahlt, ihr produziert normal.\\
\textbf{Stufe~2 eingeschränkt:} bis zur Hälfte der Kosten unbezahlt, eure Produktion halbiert sich in der \emph{nächsten} Runde (aufgerundet).\\
\textbf{Stufe~3 Stillstand:} mehr als die Hälfte unbezahlt, ihr produziert in der nächsten Runde nichts, und die Station verliert 1~Stabilität.
\end{metrobox}

\begin{metrobox}[title={4 \ Stabilität (Start: 10 Punkte)},fonttitle=\bfseries]{cM}
Die Stabilität der Station sinkt, wenn Gruppen im Stillstand stehen, die Krankheitskarte nicht bezahlt wird oder ein Ereignis ungelöst bleibt. Sind alle fünf Gruppen am Rundenende im Normalbetrieb, steigt sie um 1. Bei \textbf{0} bricht die Station auseinander: dann verlieren alle.
\end{metrobox}

\begin{metrobox}[title={5 \ Regierung, Verträge, Zwangsmaßnahme},fonttitle=\bfseries]{cW}
\textbf{Gemeinschaftsvorrat:} Ihr dürft beschließen, Marker in einen gemeinsamen Vorrat zu legen und wer ihn verteilt.\\
\textbf{Verträge} dürft ihr schriftlich festhalten. Sie sind bindend, aber nicht automatisch durchsetzbar.\\
\textbf{Eine Regierung braucht Marker:} Ein Beschluss ersetzt keine Marker. Was nicht da ist, kann nicht verteilt werden.\\
\textbf{Zwangsmaßnahme:} Die Wächter dürfen \emph{einmal} bis zu 2 Marker einer Gruppe einziehen. Mit Beschluss der Konferenz ist das legitim. Ohne Beschluss verliert die Station 2~Stabilität.
\end{metrobox}
